\documentclass[10pt,a4paper]{article}
\usepackage[utf8]{inputenc}
\usepackage[T1]{fontenc}
\usepackage[spanish,es-noquoting]{babel}
\usepackage[margin=12mm]{geometry}
\usepackage{lmodern,microtype,booktabs,array,xcolor,tikz,hyperref}
\usetikzlibrary{arrows.meta,positioning,calc}
\hypersetup{colorlinks=true,urlcolor=blue!55!black}
\pagestyle{empty}
\setlength{\parindent}{0pt}
\setlength{\parskip}{3pt}
\definecolor{navy}{HTML}{16354A}
\definecolor{teal}{HTML}{087E8B}
\definecolor{pale}{HTML}{EAF3F6}
\newcommand{\subhead}[1]{\vspace{3pt}{\color{navy}\bfseries #1}\par}
\begin{document}
{\LARGE\bfseries\color{navy} Tres agentes para normativa UdeC}\par
{\small Álvaro Contreras y Pablo Cortés \hfill Generative Artificial Intelligence (580694)}\par
{\footnotesize Deliverable 3 \hfill 30 de octubre de 2026}
\par\hrule\vspace{4pt}

\subhead{Continuidad, tarea y criterio}
{\small Preguntas en español sobre RG, RI-FI y calendario académico 2026: mismo corpus de D2 (197 unidades; 201 fragmentos) y mismas 50 preguntas de desarrollo. Se busca resolver omisiones, remisiones incompletas y citas que no respaldan la respuesta. Una salida correcta entrega lo solicitado, con respaldo pertinente y sin añadidos falsos; rechaza premisas falsas cuando corresponde. La evaluación externa usa la pauta D2; aceptar internamente no equivale a acertar.}

\subhead{Arquitectura: decisiones observables, no una secuencia fija}
\begin{center}
\begin{tikzpicture}[>=Stealth,font=\scriptsize,
agent/.style={draw=navy,fill=pale,rounded corners=2pt,minimum height=8mm,minimum width=20mm,align=center},
tool/.style={draw=teal,fill=teal!5,minimum height=7mm,minimum width=28mm,align=center},
endnode/.style={draw=gray,rounded corners,align=center,minimum height=7mm},
lab/.style={fill=white,inner sep=1pt,font=\tiny,align=center}]
\node[endnode] (q) at (0,0) {Pregunta};
\node[agent] (i) at (3.4,0) {Investigador};
\node[agent] (r) at (7.2,0) {Redactor};
\node[agent] (a) at (11,0) {Auditor};
\node[endnode] (f) at (14.2,0) {Respuesta};
\node[tool] (b) at (0.6,-2.4) {buscar\_normativa\\E5 y recuperador D2};
\node[tool] (l) at (3.7,-2.4) {leer\_articulo\\unidades completas};
\node[tool] (n) at (6.8,-2.4) {inventario\\estructura de la copia};
\node[endnode] (x) at (11.5,-2.4) {Fin sin validar\\límite de llamadas / formato / contexto};
\draw[->] (q) -- node[lab,left]{búsqueda inicial} (b);
\draw[<->] (i.south west) -- node[lab,above,sloped]{consulta / fuentes} (b.north east);
\draw[<->] (i.south) -- node[lab,right,pos=0.78]{número /\\texto} (l.north);
\draw[<->] (i.south east) -- node[lab,above,sloped]{documento /\\inventario} (n.north west);
\draw[->] (i) -- node[lab,above]{redactar:\\fuentes y aviso} (r);
\draw[->] (r) -- node[lab,above]{auditar:\\borrador y citas} (a);
\draw[->] (a) -- node[lab,above]{aceptar:\\respuesta} (f);
\draw[->] (r.north) to[out=145,in=35] node[lab,above]{investigar: carencia / control de citas o repetición} (i.north);
\draw[->] (a.north) to[out=145,in=35] node[lab,above]{corregir: objeción} (r.north);
\draw[->] (a.north east) -- ++(0,1.5) -| node[lab,pos=0.3,above]{investigar: respaldo pendiente / dos correcciones sin evidencia nueva} (i.north west);
\draw[->,dashed] (i.south) -- ++(0,-3.15) -| node[lab,pos=0.6,below]{límite alcanzado} (x.south);
\draw[->,dashed] (r.south) -- ++(0,-2.95) -| (x.south);
\draw[->,dashed] (a.south) -- (x.north);
\end{tikzpicture}
\end{center}
\vspace{-4pt}
{\footnotesize Las flechas dobles intercambian solicitudes y fuentes. Cada llamada recibe pregunta, evidencia, observaciones, borrador, mensaje y presupuesto; devuelve una decisión JSON validada. El investigador elige herramientas o entrega; el redactor redacta o pide datos; el auditor acepta, corrige o investiga. El controlador registra además transferencias por cita desconocida o repetición, e invalida borradores obsoletos.}

\subhead{Técnicas y ablaciones}
{\small
\begin{tabular}{@{}p{29mm}p{67mm}p{69mm}@{}}
\toprule
Agente & Técnica y fallo que aborda & Ablación de un bloque externo\\
\midrule
Investigador & Descomposición explícita: cobertura de partes y remisiones & \texttt{sin\_descomposicion}\\
Redactor & Lista de completitud: omisiones y condiciones & \texttt{sin\_completitud}\\
Auditor & Verificación por afirmación: citas y añadidos & \texttt{sin\_verificacion}\\
\bottomrule
\end{tabular}}\par
{\footnotesize Se retira solo el bloque de prompting indicado; se mantienen roles, herramientas, protocolo y límites. Son contribuciones incrementales: D2 ya contiene expansión de remisiones y el auditor genérico permanece.}

\subhead{Resultados empíricos: evaluación de 300 respuestas (6 condiciones $\times$ 50 preguntas)}
{\footnotesize
\begin{tabular}{@{}lrrrrrr@{}}
\toprule
Condición experimental & Total (50) & P01--P40 & Llam./q & Tok. In/q & Tok. Out/q & Latencia/q \\
\midrule
\texttt{baseline\_directo} & 1/50 (2,0\%) & 1/40 (2,5\%) & 1,0 & 119 & 65 & 6,7\,s \\
\texttt{rag\_d2} & 40/50 (80,0\%) & 30/40 (75,0\%) & 1,0 & 1.320 & 95 & 15,1\,s \\
\texttt{agentes\_completo} & 34/50 (68,0\%) & 34/40 (85,0\%) & 5,1 & 12.009 & 887 & 152,3\,s \\
\texttt{sin\_descomposicion} & 34/50 (68,0\%) & 32/40 (80,0\%) & 5,4 & 12.598 & 933 & 161,5\,s \\
\texttt{sin\_completitud} & 34/50 (68,0\%) & 33/40 (82,5\%) & 5,2 & 11.711 & 816 & 142,1\,s \\
\texttt{sin\_verificacion} & 32/50 (64,0\%) & 31/40 (77,5\%) & 4,9 & 10.515 & 867 & 139,2\,s \\
\bottomrule
\end{tabular}}\par
{\footnotesize En las 40 preguntas sustantivas de la normativa, la arquitectura multi-agente supera a D2 por +10,0 puntos porcentuales (85,0\% vs.\ 75,0\%), resolviendo 9 de las 10 fallas históricas de D2 (P01, P11, P13, P15, P24, P25, P28, P30, P38). En P41--P50 (artículos inexistentes), el auditor rechaza convalidar citas sin respaldo y agota presupuesto en abstención técnica, evitando alucinaciones pero cediendo frente a la abstención directa de D2. Las ablaciones confirman el impacto de cada técnica: retirar verificación reduce el acierto global en 4,0 pp (32/50), mientras que prescindir de descomposición o completitud induce fallas críticas por desalineación de citas y límite de salidas inválidas (P06, P07, P12, P28). El coste asciende a 5,12 llamadas y 152,3\,s por caso.}

\subhead{Reproducción y límites}
{\small Los tres agentes comparten Qwen3-4B abierto, NF4 en T4, sin thinking, temperatura 0,7; E5 solo recupera. Por entrada: 12 llamadas, 6 herramientas (incluida búsqueda inicial), 3 salidas inválidas; contexto 8.192, salida 1.024 y margen 64 tokens. Baseline y D2 conservan sus prompts originales. Cada ZIP congela código, corpus, prompts, semillas, respuestas, tiempos y trazas; los juicios se vinculan por hashes. Las 50 preguntas se usaron en desarrollo: no demuestran generalización.}\par
{\footnotesize \textbf{Repositorio:} \href{https://github.com/Pacortes2021/GENERATIVE-ARTIFICIAL-INTELLIGENCE-DELIVERABLE-1/tree/main/Deliverable3_Agentes}{github.com/Pacortes2021/GENERATIVE-ARTIFICIAL-INTELLIGENCE-DELIVERABLE-1}\\
\textbf{Video demostración ($\le$ 3 min):} enlace a grabación pública según \texttt{guion\_video.md}.}
\end{document}